% Диаграмма развёртывания системы автоматизации приёмной комиссии.
% Чёрно-белое оформление по конвенции БНТУ.
% Сборка: xelatex deployment.tex (см. build-diagrams.sh).

\documentclass[border=4mm]{standalone}

\usepackage{fontspec}
\setmainfont{Times New Roman}

\usepackage{tikz}
\usetikzlibrary{positioning,calc,arrows.meta,fit,shapes.geometric}

\begin{document}
\begin{tikzpicture}[
  font=\fontsize{11pt}{13pt}\selectfont,
  every node/.style={align=center},
  container/.style={
    draw=black, thick, rectangle,
  },
  containerlabel/.style={
    anchor=north west,
    font=\fontsize{11pt}{13pt}\selectfont\bfseries,
    inner sep=1pt, fill=white,
  },
  comp/.style={
    draw=black, thick, fill=white, rectangle,
    minimum height=14mm, minimum width=42mm, inner sep=2mm,
  },
  art/.style={
    draw=black, thick, fill=white, rectangle,
    minimum height=11mm, minimum width=42mm, inner sep=2mm,
  },
  cyl/.style={
    draw=black, thick, fill=white,
    cylinder, shape border rotate=90, aspect=0.25,
    minimum width=44mm, minimum height=17mm,
    inner sep=2mm,
  },
  arrow/.style={-{Stealth[length=2.5mm]}, thick},
  dotlink/.style={dashed, thin},
  edgelbl/.style={
    font=\fontsize{9.5pt}{11pt}\selectfont,
    align=center, fill=white, inner sep=2pt,
  },
]

% ============================================================
% РАБОЧАЯ СТАНЦИЯ (слева)
% ============================================================
\node[comp] (browser) {Веб-браузер\\(Chrome, Firefox, Safari)};
\node[art, below=18mm of browser] (token) {JWT-токен\\в localStorage};
\draw[dotlink] (browser) -- (token);

\node[container, fit=(browser)(token), inner xsep=10mm, inner ysep=10mm] (ws) {};
\node[containerlabel] at ([xshift=2mm, yshift=-1mm]ws.north west)
  {Рабочая станция пользователя};

% ============================================================
% СЕРВЕР ПРИЛОЖЕНИЙ (справа) — сначала строим внутренности
% ============================================================
% Верхний ряд: nginx и backend.
\node[comp, right=72mm of browser] (nginx)
  {Веб-сервер\\статических файлов\\(nginx)};
\node[comp, right=22mm of nginx] (backend)
  {Серверная часть\\(ASP.NET Core 8,\\Kestrel)};

% Артефакт SPA — под nginx.
\node[art, below=12mm of nginx] (spa) {Сборка SPA\\(HTML, JS, CSS)};
\draw[dotlink] (nginx) -- (spa);

% Контейнер liquibase — ниже, слева внутри docker.
\node[comp, below=32mm of spa] (liquibase)
  {Liquibase 5.0\\(миграции схемы)};
\node[container, fit=(liquibase), inner xsep=5mm, inner ysep=14mm] (liqbox) {};
\node[containerlabel, align=left] at ([xshift=2mm, yshift=-1mm]liqbox.north west)
  {Контейнер\\liquibase};

% Контейнер с postgres и томом — справа внутри docker.
\node[cyl, right=28mm of liquibase] (postgres) {PostgreSQL 18\\alpine};
\node[art, below=12mm of postgres] (volume) {Том\\postgres\_data};
\draw[dotlink] (postgres) -- (volume);
\node[container, fit=(postgres)(volume), inner xsep=5mm, inner ysep=14mm] (pgbox) {};
\node[containerlabel, align=left] at ([xshift=2mm, yshift=-1mm]pgbox.north west)
  {Контейнер\\bntuapplicants-postgres};

% Обёртка Docker Engine.
\node[container, fit=(liqbox)(pgbox), inner xsep=7mm, inner ysep=10mm] (docker) {};
\node[containerlabel] at ([xshift=2mm, yshift=-1mm]docker.north west)
  {Docker Engine};

% Обёртка всего сервера.
\node[container, fit=(nginx)(backend)(spa)(docker),
      inner xsep=8mm, inner ysep=11mm] (server) {};
\node[containerlabel] at ([xshift=2mm, yshift=-1mm]server.north west)
  {Сервер приложений};

% ============================================================
% СТРЕЛКИ
% ============================================================
% Обе стрелки из browser выходят справа симметрично относительно центра
% блока: REST/JSON — выше центра на 4mm (уходит вверх), статика — ниже
% центра на 4mm (идёт прямо к nginx).

% browser -> nginx (статика). Прямая горизонтальная, ниже центра.
\coordinate (br_nginx_a) at ([yshift=-4mm]browser.east);
\coordinate (br_nginx_b) at ([yshift=-4mm]nginx.west);
\draw[arrow] (br_nginx_a) -- (br_nginx_b);
\node[edgelbl] at ($(br_nginx_a)!0.5!(br_nginx_b)$)
  {HTTPS\\(статика)};

% browser -> backend (REST/JSON). Маршрут «вправо -> вверх -> вправо -> вниз»:
% выходим горизонтально из browser выше центра, пересекая правую границу ws,
% затем поднимаемся выше уровня сервера, идём вправо до точки над backend
% и спускаемся в backend.north. Метку размещаем на верхнем горизонтальном
% участке выше всех контейнеров.
\coordinate (br_back_start) at ([yshift=4mm]browser.east);
\coordinate (br_back_a) at ([xshift=4mm]ws.east |- br_back_start);
\coordinate (br_back_b) at ([yshift=20mm]br_back_a);
\coordinate (br_back_c) at (br_back_b -| backend.north);
\draw[arrow] (br_back_start) -- (br_back_a) -- (br_back_b) -- (br_back_c) -- (backend.north);
\node[edgelbl] at ($(br_back_b)!0.5!(br_back_c)$)
  {HTTPS, REST/JSON\\(порт 5059)};

% backend -> postgres (SQL). Маршрут: вправо от backend, вниз вдоль правой
% стороны server, затем влево внутрь pgbox до postgres.east — обходит заголовок
% контейнера bntuapplicants-postgres.
\coordinate (bp_a) at ([xshift=12mm]backend.east);
\coordinate (bp_b) at (bp_a |- postgres.east);
\draw[arrow] (backend.east) -- (bp_a) -- (bp_b) -- (postgres.east);
\node[edgelbl] at ([yshift=-20mm]bp_a) {SQL\\(порт 5432)};

% liquibase -> postgres (SQL при обновлении). Горизонтальная стрелка.
\draw[arrow] (liquibase.east) -- (postgres.west)
  node[edgelbl, midway, above=2mm]
  {SQL (порт 5432,\\при обновлении)};

\end{tikzpicture}
\end{document}
